% bank_fintech_locus/drafts/en/manuscript.tex
% P4 manuscript skeleton — M0 G0 scaffolded per B-route plan v2.55 (2026-05-17)
% M1 enrichment (2026-07-03): intro/lit/hypotheses/methods/data/results/conclusion
%   + validity-commitment appendix + network-roles appendix
% Compile: xelatex + biber + xelatex × 2

% papers_assets/_shared/preamble_en.tex
% Shared English LaTeX preamble for all 5 paper drafts.
% P4.1 (plan-r3 derivative-paper extension, 2026-05-15)
% Compile: xelatex + biber + xelatex × 2

\documentclass[12pt,a4paper]{article}

% --- Fonts & Encoding ---
\usepackage{fontspec}
\setmainfont{Times New Roman}[Ligatures=TeX]

% --- Math ---
\usepackage{amsmath, amssymb, amsthm}

% --- Tables & Figures ---
\usepackage{booktabs}
\usepackage{graphicx}
\usepackage{longtable}
\usepackage{rotating}
\usepackage{caption}
\captionsetup{font=small,labelfont=bf,skip=4pt}

% --- Layout ---
\usepackage[margin=1in]{geometry}
\usepackage{setspace}
\onehalfspacing
\usepackage{hyperref}
\hypersetup{colorlinks=true,linkcolor=blue,citecolor=blue,urlcolor=blue}

% --- Bibliography (biblatex + biber) ---
\usepackage[backend=biber,style=apa,natbib=true]{biblatex}

% --- Useful macros ---
\newcommand{\code}[1]{\texttt{#1}}
\newcommand{\NoteThis}[1]{\textbf{\textcolor{red}{[TODO: #1]}}}
\usepackage[dvipsnames]{xcolor}

% --- Theorem environments ---
\newtheorem{prop}{Proposition}
\newtheorem{hyp}{Hypothesis}


\addbibresource{refs.bib}

\title{Bank Fintech Locus:\\Layer (Underlying Technology) vs Scenario (Business Function) Market Reactions across Three Waves}
\author{\NoteThis{Author names}}
\date{\today}

\begin{document}
\maketitle

\begin{abstract}
\NoteThis{Final abstract pending — target $\approx$ 200 English words; headline numbers
await production-corpus (T=19) regressions.}
Following the Digital Finance Cube of Gomber, Koch, and Siering (2017), we classify
Chinese banks' fintech media discourse into two locus categories — underlying technology
(blockchain / AI / cloud) versus business-function scenario (payments / lending / wealth
management) — and test a structurally paired set of theoretical predictions:
technology-layer discourse should map into technological innovation (citation-weighted
patents, in the Kogan et al.\ 2017 lineage; citation-count weighting for unlisted banks),
whereas scenario-layer discourse should map into product innovation and fee-based
business (noninterest income share). We build a bank--year panel from a production corpus
of 18{,}198{,}978 Chinese media reports (the full census basis, \emph{before} the
bank--fintech co-occurrence hard filter; 2007--2025, 110 banks,
$T=19$); article-level locus labels
come from a majority vote across three independent-vendor LLMs. Identification combines
bank and year fixed effects, layer $\times$ event studies around four policy shocks (the
Yu'ebao launch, the Ant IPO suspension, and two data-export regulations), and a
leave-one-out peer-locus instrument within bank-type--year cells. Event-window CARs for
the listed-bank subsample serve only as an extension. We find \NoteThis{key result}\@.
\end{abstract}

\section{Introduction}\label{sec:intro}

Over the past decade and a half, Chinese banking has undergone a digital transformation
of rare scale: from the 2013 Yu'ebao launch that ignited the internet-finance era (wave
1, 2013--2015), through the blockchain and early-AI technology race (wave 2, 2016--2018),
to the large-language-model wave since 2022 (wave 3, 2022--2024), with a regulatory
consolidation period (2019--2021) following the suspension of Ant Group's IPO in November
2020. Throughout, the media has been the most important information intermediary for
banks' technology strategies: what a bank talks about — underlying technologies such as
blockchain, AI, and cloud computing, or application scenarios such as payments, lending,
and wealth management — constitutes a high-frequency public signal of its innovation
locus. Yet most existing work treats bank fintech activity as a homogeneous scalar,
ignoring the orthogonal distinction between a technology dimension and a
business-function dimension long established in the digital-finance taxonomy
\parencite{gomber2017cube,gomber2018fintech}.

The core question of this paper is whether the locus of a bank's fintech media discourse
— dominated by underlying technology or by application scenarios — differentially
predicts the \emph{form} of its innovation output. Following the Digital Finance Cube of
\textcite{gomber2017cube}, we aggregate article-level locus labels into two mutually
exclusive dominance indicators at the bank--year level (technology-layer dominance,
\code{locus\_layer\_full}, and scenario dominance, \code{locus\_scenario\_full}) and
derive structurally paired predictions: technology-layer investment alters a bank's
information production function and should yield patentable technological innovation
(citation-weighted patents), whereas scenario-layer investment alters distribution and
product interfaces and should yield fee-based business (noninterest income share)
\parencite{yoo2010layered,boot2021fintech}.

On the data side, we start from a raw census of 18{,}198{,}978 media reports (2007--2025,
3{,}740 outlets) and take this full census as the production target corpus ---
18{,}198{,}978 media reports covering 110 banks over
2007--2025 ($T=19$) — a basis that sits \emph{before} the bank--fintech co-occurrence
hard filter, which, together with the full-track retention gate, substantially shrinks
the set of articles entering the analysis. Article-level locus is classified by a
majority vote across three independent-vendor large language models (temperature $=0$,
byte-frozen prompt), under a full offline-cache reproducibility protocol
\parencite{ludwig2025llmecon,carlson2026llmannotation}. On the identification side, the
baseline uses bank and year fixed effects with bank-level clustering, complemented by
layer $\times$ event studies around four policy shocks (E1 Yu'ebao launch, 2013-06-13;
E4 Ant IPO suspension, 2020-11-03; E5a/E5b data-export regulations, 2022-09-01 /
2023-09-30) estimated with heterogeneity-robust estimators
\parencite{sunabraham2021eventstudy,borusyak2024revisiting,callaway2021difference}, and
a leave-one-out peer-locus mean instrument within bank-type--year cells
\parencite{manski1993reflection,leary2014peereffects}. Our inference discipline
pre-declares a single confirmatory hypothesis, labels everything else exploratory, and
controls the family-wise error rate with Romano--Wolf stepdown
\parencite{romanowolf2005stepdown}.

The paper makes four contributions. First, measurement: we operationalize the Digital
Finance Cube's technology $\times$ business-function taxonomy as a bank--year locus panel
measure on a large media corpus, with fully disclosed LLM-annotation validity commitments
and a reproducibility protocol. Second, design: the binary-treatment-to-binary-outcome
structural pairing distinguishes this paper from sibling studies on the same corpus
family that measure continuous tone, exposure intensity (attention), the
substantive-versus-symbolic dimension, or the in-house-versus-open collaboration
dimension — locus measures the technology-layer attribute of the announcement itself,
not its sentiment, salience, or organizational source. Third, identification: within an
annual panel we combine leave-one-out peer instruments, partial-identification
sensitivity for event studies \parencite{rambachan2023honestdid}, and a multi-ladder
robustness architecture in which every ladder enters its own just-identified
specification. Fourth, context: we bring systematic evidence from China — the world's
largest digital-finance market, including many unlisted banks — to a bank-technology
literature dominated by U.S. patents and listed firms
\parencite{chen2019fintechvalue,lerner2024finpatents}.

We preview the main findings as follows: \NoteThis{results preview — direction and
magnitude for H1/H2/H3/H4, to be backfilled from production regressions}.

The remainder proceeds as follows. Section \ref{sec:lit} reviews four strands of
literature; Section \ref{sec:hypo} develops the hypotheses; Section \ref{sec:method}
presents the sample, locus measurement, and identification strategy; Section
\ref{sec:data} describes data sources; Section \ref{sec:results} reports results; and
Section \ref{sec:conclusion} concludes. Appendices contain the measurement-validity
commitments and the exploratory network-roles analysis.

\section{Literature Review}\label{sec:lit}

\subsection{Fintech taxonomy and the layered architecture of digital innovation}

The direct conceptual source of our locus construct is the Digital Finance Cube of
\textcite{gomber2017cube}: digital-finance innovation can be characterized along three
orthogonal dimensions — technologies, business functions, and institutions — of which
the first two are the theoretical prototype of our underlying-technology versus
application-scenario dichotomy; \textcite{gomber2018fintech} provide the functional map
of the fintech revolution. Information-systems theory supplies the architectural
foundation: \textcite{yoo2010layered} show that digital products consist of device,
network, service, and contents layers, and that where the locus of control sits
determines the organizing logic of innovation. In finance,
\textcite{boot2021fintech} decompose fintech into an information axis and a
communication axis, isomorphic to our layer (information production) / scenario
(distribution and communication) split; \textcite{vives2019digital} surveys incumbent
banks' strategy space under digital disruption; \textcite{parlour2022payments} and
\textcite{berg2022fintechlending} model competition at the payments layer and the
lending application layer, respectively; the underlying-technology character of
blockchain and tokenomics is analyzed in
\parencite{conghe2019blockchain,cong2021tokenomics}. Finally,
\textcite{adner2010ecosystem} show that upstream (component) versus downstream
(complement) technological interdependence affects firm performance through
fundamentally different channels — the layer-versus-scenario split is thus an
economically grounded distinction, not an arbitrary label.

\subsection{Bank technology adoption and innovation output}

A second strand links banks' technology investment to measurable innovation output.
\textcite{kogan2017technological} establish the classic lineage of valuing patents by
market reactions; \textcite{chen2019fintechvalue} use machine learning to classify
patent text into fintech technology categories and show that private value differs
sharply across categories — the direct top-journal precedent for ``different technology
loci carry different economic value''; \textcite{lerner2024finpatents} characterize who
produces financial-innovation patents. On the IT-input side, \textcite{he2026lendingit}
decompose bank IT spending into distinct technology loci; \textcite{vivesye2025itcompetition}
provide a theory of information technology and lender competition;
\textcite{pierri2022techbanking} show that technology adoption shaped lending
performance in a crisis; and \textcite{kwan2024digitalcapability} offer a bank-level
digital-capability index as a measurement precedent. External validity in the Chinese
context comes from \textcite{hau2024fintechcredit} (fintech credit evidence based on Ant
Group data) and \textcite{alfhaili2026fintechequity} (banks' fintech equity investment
and bank innovation); on big data in finance generally, see
\parencite{goldstein2021bigdata,goldstein2019fintech}.

\subsection{LLM-based text measures as econometric variables}

A third strand concerns the measurement method itself. \textcite{ludwig2025llmecon}
develop an applied econometric framework for LLM-generated variables entering downstream
analysis, distinguishing prediction from estimation uses and emphasizing how annotation
error propagates into inference; \textcite{carlson2026llmannotation} provide
foundational guidelines for LLM annotation in management research, warning about
systematic bias and prompt sensitivity. Our design responds directly to these warnings:
a majority vote across three independent vendors hedges single-vendor systematic bias;
a byte-frozen prompt and zero temperature remove stochasticity; an offline cache
guarantees byte-level reproducibility; and Appendix \ref{app:validity} reports the
pre-committed expert-annotation validity checks.

\subsection{Identification and inference}

A fourth strand underpins the identification design. The identification roots of peer
instruments are \textcite{manski1993reflection}'s reflection problem and
\textcite{leary2014peereffects}'s leave-one-out peer construction in finance; the two
identification routes for generalized Bartik-type instruments are given by
\textcite{goldsmithpinkham2020bartik} (exogenous shares) and
\textcite{borusyak2022shiftshare} (exogenous shocks). Event-study estimation under
staggered treatment and heterogeneous effects follows
\parencite{sunabraham2021eventstudy,callaway2021difference,borusyak2024revisiting};
partial-identification sensitivity to local parallel-trends violations follows
\textcite{rambachan2023honestdid}; family-wise error control follows
\textcite{romanowolf2005stepdown}; and few-cluster inference follows
\textcite{roodman2019boottest}'s wild cluster bootstrap.

\subsection{Positioning}

In sum, the taxonomy literature establishes the layer-versus-scenario distinction but
lacks large-sample bank-panel evidence; the bank-technology literature measures IT
spending or patent output but does not directly observe the technology positioning of
banks' public discourse; and the LLM-measurement literature provides tools but few
applications embedded in a full identification framework. This paper sits at the
intersection: it converts a taxonomic construct into a panel variable using a
discipline-constrained LLM measure and tests its structurally paired predictions about
the form of innovation output with modern causal-inference tools.

\section{Theory and Hypotheses}\label{sec:hypo}

Our theoretical skeleton is that the locus of a bank's fintech media discourse reflects
where its digital investment sits in the layered architecture
\parencite{gomber2017cube,yoo2010layered}, and that investment at different layers maps
into different forms of innovation output through different economic mechanisms
\parencite{boot2021fintech,adner2010ecosystem}. All hypotheses below are theoretical
predictions, not results; only H1 is the pre-declared confirmatory hypothesis, while
H2--H4 are exploratory.

\paragraph{H1 (confirmatory primary hypothesis: technology layer $\to$ technological
innovation).}
Discourse about underlying technologies (blockchain / AI / cloud) reflects a bank's
genuine investment orientation toward general-purpose technologies. Such investment
alters the bank's information production function, and its output naturally crystallizes
as patentable technological solutions: \textcite{chen2019fintechvalue} show that patents
in different fintech technology categories carry significant and heterogeneous private
value, and \textcite{he2026lendingit} show that competitive pressure induces banks to
scale up software-type IT investment. \emph{Theoretical prediction}: the
technology-layer dominance indicator \code{locus\_layer\_full} (derived from
\code{tech\_frac} $>$ \code{app\_frac}) carries a positive coefficient on
citation-weighted patent output (\code{citation\_weighted\_patent}), with bank and year
fixed effects and bank-level clustering.

\paragraph{H2 (exploratory paired hypothesis: application scenario $\to$ fee business).}
Payments, lending, and wealth management are business-function-layer innovations; their
investment changes distribution channels and product interfaces rather than the
underlying technology stack \parencite{parlour2022payments,boot2021fintech}, and their
output takes the form of fee and commission business rather than patents.
\emph{Theoretical prediction}: the scenario dominance indicator
\code{locus\_scenario\_full} carries a positive coefficient on noninterest income share
(\code{noninterest\_income\_share}); at the same time its effect on citation-weighted
patents should be markedly weaker than H1 — the placebo side of the structural pairing.
If scenario discourse predicted patents just as strongly, the locus distinction would
lack discriminant validity.

\paragraph{H3 (exploratory mediation: IT capex and R\&D intensity channels).}
If technology-layer discourse reflects real investment, its association with patent
output should partly run through observable input variables: IT capital expenditure
(\code{capex\_it}, primary mediator) and R\&D intensity (\code{r\_and\_d\_intensity},
secondary mediator). \emph{Theoretical prediction}: the indirect path
\code{locus\_layer\_full} $\to$ mediator $\to$ \code{citation\_weighted\_patent} is
positive. We estimate a Baron--Kenny decomposition \parencite{baron1986moderator} and
report the sensitivity parameter $\rho$ for violations of sequential ignorability
following \textcite{imai2010general}; because the mediators are endogenous, this
analysis is explicitly labeled an association decomposition, not causal mediation.

\paragraph{H4 (exploratory heterogeneity: three waves).}
The strength of locus effects should vary with the technological content of each wave:
wave 2 (2016--2018, blockchain / early AI) and wave 3 (2022--2024, large models) are
dominated by underlying technology, whereas wave 1 (2013--2015, internet finance) is
dominated by payments and wealth-management scenarios. \emph{Theoretical prediction}:
the effect of \code{locus\_layer\_full} on citation-weighted patents is stronger in
wave\_2 and wave\_3; the effect of \code{locus\_scenario\_full} on noninterest income
share is stronger in wave\_1; the consolidation period (2019--2021) serves as the
comparison period (three waves of exactly three years each, strictly symmetric).

\section{Methodology}\label{sec:method}

\subsection{Sample and Data}

$N \times T \approx 3000\text{--}9000$ bank--years, 2007--2025 (per spec.sample).
% TODO: 3000--9000 区间系旧口径 T=18 下的推导值, 待全量重跑后按 T=19 更新。
\NoteThis{paper\_spec.sample (n\_bank\_min=200 / obs\_target\_min=3000) is a
bank-universe expansion target and conflicts with the current production count of 110
banks $\times$ $T=19$ (at most 2090 bank--years); sample sizes to be reconciled and
backfilled after sample-frame enrichment and production-panel finalization.}

Corpus construction follows two explicitly separated accounting bases. The \emph{raw
census basis} (before filtering and deduplication; for the data appendix): 183
deduplicated source files containing 18{,}198{,}978 media reports spanning 2007-01-01 to
2025-12-31 across 3{,}740 outlets; newspaper-type reports account for 92.93\% of all
rows (16{,}912{,}361 rows) and are retained, while wire-service, magazine, government,
and web types are dropped row-wise; exactly unique records number 18{,}177{,}370
(duplication rate 0.12\%). The \emph{production target basis} (the full census scope,
\emph{before} the bank--fintech co-occurrence hard filter):
18{,}198{,}978 media reports, 2007--2025, $T=19$, 110 banks. The retained universe entering the
regressions (after the co-occurrence hard filter $+$ the \code{full}-track retention
gate) is far smaller than this basis; its size \NoteThis{to be backfilled once the
production pipeline has run}. The pre-filter basis and the post-filter regression universe must not be conflated.

The filtering chain comprises target-bank entity recognition (an asymmetric criterion
with negation and short-name filters), fintech lexicon matching, and — critically — a
co-occurrence hard filter that retains only articles in which the target bank and at
least one fintech term co-occur within a $\pm 1$-sentence window
(\code{sentence\_radius}$=1$), with the co-occurrence test covering both headline and
body (\code{match\_in\_title}$=$\code{true}). Each retained article is then materialized
into three parallel text tracks: \code{full} (headline $+$ full body), \code{ctx}
(headline $+$ entity context window), and \code{ss} (headline $+$ strict same-sentence
bank--fintech window $\pm 1$ sentence). The three tracks are mutually robust with no
ordering of primacy.

\subsection{Locus Measurement (article-level annotation at pipeline stage 21 $\to$ panel aggregation at stage 40)}

\paragraph{F1: LLM ensemble annotation.}
Article-level annotation is performed by an ensemble of large language models from three
independent vendors — \code{gpt-5.4-nano}, \code{gemini-3.1-flash-lite}, and
\code{deepseek-v4-flash} — whose outputs are combined by majority vote into a consensus
label. All calls are made with temperature $=0.0$ and seed $=42$ passed explicitly; the
prompt (v8, v8-2026-05-17) is byte-frozen. Every request--response pair is written to an
offline cache; with zero cache misses, the entire annotation can be reproduced
byte-for-byte fully offline at zero API cost
\parencite{ludwig2025llmecon,carlson2026llmannotation}. Chinese segmentation is anchored
at \code{jieba==0.42.1}.

\paragraph{F2: three text tracks.}
Each LLM dimension is annotated independently on the \code{full} / \code{ctx} /
\code{ss} tracks, yielding three track versions of \code{innovation\_locus}. The three
tracks are mutually robust with no ordering of primacy; each track enters its own
just-identified specification and is never pooled into one over-identified system (see
Section \ref{subsec:robust}).

\paragraph{F3: article-retention gate.}
Whether an article is judged a bank--fintech report and enters the analysis universe is
determined solely by the full-track verification flag
\code{llm\_verified\_full}$=$\code{true} (full-only gating); the \code{ctx} and
\code{ss} tracks are conditional-on-full robustness measurement tracks whose row
universe is a subset of the full-track retention set.

\paragraph{F4: locus labels and panel aggregation.}
The article-level locus label is a HYBRID construction: LLM majority voting as the
backbone, complemented by keyword-based confidence adjustment and an empty-vote
fallback; its validity numbers therefore inherit a keyword circularity and are not
directly comparable to pure-LLM dimensions (see Appendix \ref{app:validity}). Bank--year
aggregation (pipeline stage 40) yields the technology-layer article share
\code{tech\_frac} and the scenario article share \code{app\_frac}. The treatment pair
(\code{locus\_layer\_full}, \code{locus\_scenario\_full}) is derived in the replication
package's Stata \code{01\_data\_prep.do} (sole producer; never written to parquet):
\code{locus\_layer\_full}$=\mathbf{1}[\code{tech\_frac}>\code{app\_frac}]$ and
\code{locus\_scenario\_full}$=\mathbf{1}[\code{app\_frac}>\code{tech\_frac}]$, coded
missing if either share is missing; ties (including double zeros) set both dummies to 0,
forming the reference group; the secondary measure
\code{locus\_gap}$=$\code{locus\_layer\_full}$-$\code{locus\_scenario\_full}
$\in\{-1,0,+1\}$.

\subsection{Identification Strategy}

\paragraph{Baseline two-way fixed effects.}
The baseline specification is
\begin{equation}\label{eq:baseline}
Y_{i,t} = \beta_1\,\mathrm{Layer}_{i,t} + \beta_2\,\mathrm{Scenario}_{i,t}
 + \mathbf{X}_{i,t}'\boldsymbol{\gamma} + \alpha_i + \lambda_t + \varepsilon_{i,t},
\end{equation}
where $Y_{i,t}$ is bank $i$'s citation-weighted patents (primary outcome) or noninterest
income share (paired outcome) in year $t$; $\mathrm{Layer}$ and $\mathrm{Scenario}$ are
the binary treatment pair derived above; $\alpha_i$ and $\lambda_t$ are bank and year
fixed effects; and $\mathbf{X}_{i,t}$ collects bank-level and year-level controls (see
Section \ref{sec:data} for their current enrichment status). The main specification
clusters standard errors at the bank level (within-bank serial correlation in a $T=19$
panel); the clustering ladder and few-cluster corrections appear in the robustness
subsection.

\paragraph{Layer $\times$ event studies.}
Four policy events (E1 Yu'ebao launch, 2013-06-13; E4 Ant IPO suspension, 2020-11-03;
E5a Data-Export Security Assessment Measures, 2022-09-01; E5b Data-Export Management
Measures, 2023-09-30) provide exogenous shocks to the fintech environment. We estimate
\begin{equation}\label{eq:eventstudy}
Y_{i,t} = \sum_{k \neq -1} \delta_k\, \mathbf{1}[t - t_e = k] \times
\mathrm{Layer}^{\mathrm{pre}}_{i} + \alpha_i + \lambda_t + \varepsilon_{i,t},
\end{equation}
where $t_e$ is the event year, $\mathrm{Layer}^{\mathrm{pre}}_i$ is pre-event
technology-layer exposure, and $k=-1$ is the omitted base period. Under staggered and
heterogeneous effects we report the \textcite{sunabraham2021eventstudy}
interaction-weighted estimator and the \textcite{borusyak2024revisiting} imputation
estimator side by side, together with a joint pre-trend test (leads $-3/-2/-1=0$) and a
\textcite{rambachan2023honestdid} relative-magnitudes partial-identification sensitivity
band (run on E1, the longest window; $M$ grid $\{0,0.5,1,1.5,2\}$).

\paragraph{Leave-one-out peer instrument.}
Locus discourse may be co-determined with unobserved innovation endowments. We
instrument with the leave-one-out peer mean of \code{tech\_frac} within bank-type
$\times$ year cells:
\begin{equation}\label{eq:iv}
\code{tech\_frac}_{i,t} = \pi\, \overline{\code{tech\_frac}}^{\,\mathrm{LOO}}_{-i,g(i),t}
 + \mathbf{X}_{i,t}'\boldsymbol{\phi} + \alpha_i + \lambda_t + u_{i,t},
\end{equation}
where $g(i)$ is bank $i$'s type group. The identification logic follows the peer-effects
literature: peers' discourse orientation shifts the focal bank's strategic language
choice but, conditional on two-way fixed effects and controls, does not directly enter
its patent production function
\parencite{manski1993reflection,leary2014peereffects}. The primary instrument is the
contemporaneous LOO mean (\code{leave\_one\_out\_peer\_locus\_iv}); the fallback is its
one-year lag (\code{leave\_one\_out\_peer\_locus\_iv\_lag1}, sidestepping contemporaneous
reflection); the \code{ctx} / \code{ss} tracks receive their own same-track LOO
instruments (\code{lpeer\_locus\_iv\_ctx\{,\_lag1\}} /
\code{lpeer\_locus\_iv\_ss\{,\_lag1\}}). Each endogenous-variable--instrument pair
enters its own just-identified 2SLS specification; we never stack multiple instruments
into one over-identified system, and hence report no Hansen J — consistency across
variants is itself the robustness evidence. In weak-instrument scenarios
(Kleibergen--Paap F below conventional thresholds) we read Anderson--Rubin confidence
sets instead. At the project level, generalized Bartik and media-outlet house-style
instrument engines are available as extensions
\parencite{goldsmithpinkham2020bartik,borusyak2022shiftshare}.

\paragraph{Inference discipline.}
The single pre-declared confirmatory hypothesis is H1 (\code{locus\_layer\_full} $\to$
\code{citation\_weighted\_patent}, two-way fixed effects); all remaining specifications
(H2--H4, track twins, cross-vendor twins, secondary outcomes, heterogeneity splits,
network roles) are explicitly labeled exploratory. The family-wise error rate of the
$\{\code{locus\_layer\_full}, \code{locus\_scenario\_full}\}$ treatment family is
controlled by Romano--Wolf stepdown with 1{,}000 replications
\parencite{romanowolf2005stepdown}; with few clusters or few-treated event studies we
fall back on the wild cluster bootstrap with 9{,}999 replications
\parencite{roodman2019boottest}.

\subsection{Embedding-Based Measures (Co-Primary Dual Track; F5)}

All embedding-derived features are computed under two co-primary
embedding models from independent vendors, reported side by side as
mutual robustness: Alibaba \code{Qwen3-Embedding-8B} (4{,}096-d,
accessed via the SiliconFlow API; canonical unsuffixed columns) and
OpenAI \code{text-embedding-3-large} (3{,}072-d; \code{\_lg}-suffixed
columns), under strict column parity (each feature exists in both
tracks or in neither). Each track enters its own just-identified
specification; the two are never pooled into one over-identified
system. Earlier pipeline vintages used a Google Gemini embedding in
the canonical slot; as of pipeline v3.15.0 (June 2026) the canonical
slot is Alibaba \code{Qwen3-Embedding-8B}, with all embedding-derived
features regenerated under a controlled full re-embedding and
embedding-space semantic validity gated by a Spearman
rank-correlation check ($\rho \geq 0.80$). The sentiment-annotation
LLM ensemble (F1) is unaffected by this change: embedding models and
sentiment LLMs are disjoint pipeline components.

We stress that the locus treatment itself is an LLM classification, not an
embedding-derived measure; the dual-embedding ladder applies to embedding-derived
auxiliary measures and control features.

\subsection{Robustness Design}\label{subsec:robust}

Our robustness architecture consists of several non-crossing ladders under one unified
discipline: within each ladder, all variants are \emph{mutually robust with no ordering
of primacy}, each enters its \emph{own just-identified} specification, and variants of
the same construct are \emph{never} pooled into one over-identified Hansen J system.

\paragraph{Three-track alternative classifiers with per-track IV mirrors.}
The same locus construct is recomputed under the \code{full} / \code{ctx} / \code{ss}
text bases, yielding the ctx twin (\code{tech\_frac\_lc}) and the ss twin
(\code{tech\_frac\_ls}) of \code{tech\_frac}; each track's endogenous variable is paired
with its same-track LOO instrument, just-identified track by track, and full $\times$
ctx $\times$ ss must never be composed into an over-identified system. Row-universe
discipline: the retention gate is full-only; ctx / ss are conditional-on-full.

\paragraph{Cross-vendor per-model classification ladder.}
Cross-vendor robustness of the treatment: the individually classified
$\{\code{gpt54nano},\code{gemini31},\code{deepseekv4}\}$\code{\_tech\_frac\_full}
variants stand side by side with the consensus-vote version as co-equal mutual
robustness arms; each per-model variant's denominator is that model's own count of
verified articles. The consensus label is not a bare majority vote (it includes the
verification gate, vote cleaning, and a keyword fallback), so the per-model ladder is a
genuinely independent measurement arm. True family-wise error control for the
cross-model family is carried by the project's root-level Romano--Wolf template
(Table A12, Panel C); each vendor is just-identified on its own, with no pooled Hansen J.

\paragraph{IV variant ladder and weak-instrument inference.}
The primary instrument, its lagged fallback, and the per-track mirrors form six
endogenous-variable--instrument pairs reported one by one, each just-identified (see
the identification subsection); when the KP F is low we read Anderson--Rubin confidence
sets; no over-identification test is ever constructed.

\paragraph{Clustering ladder and wild bootstrap.}
Beyond bank-level clustering in the main specification, we report bank-only, bank
$\times$ year two-way, and bank-type $\times$ year two-way clustering side by side; with
$G<50$ clusters or few-treated event studies we recheck with the boottest wild cluster
bootstrap (9{,}999 replications) \parencite{roodman2019boottest}.

\paragraph{Honest-DiD parallel-trends sensitivity.}
Sensitivity of event-study conclusions to local parallel-trends violations is presented
as \textcite{rambachan2023honestdid} relative-magnitudes partial-identification bands
(E1, $M \in \{0,0.5,1,1.5,2\}$), reported alongside the heterogeneity-robust estimators.

\paragraph{Cell decomposition and clarity robustness table.}
A $2^2$ cross-construct cell decomposition of locus against the adjacent collab\_mode
construct (from the sibling collaboration paper) establishes dimensional independence
(Table 6); collab\_mode serves only as an alternative-specification control in Table 5,
not a parallel treatment. The language-clarity robustness table (Table A7) is reported
dual-track. The cancelled title-track table does not reappear.

\paragraph{Scope statement.}
Other measurement ladders in the project's corpus family (political-word weights and
linguistic-rule ablations for dictionary tone, rare-word frequency cutoffs, the
forward-looking-statement two-arm design, BERT per-model sentiment families, etc.)
belong to sibling papers' construct families, are unrelated to this paper's
LLM-classified locus, and do not enter this paper's robustness architecture.

\section{Data}\label{sec:data}

\paragraph{Media corpus.}
The text-side data are the bank--fintech media corpus described in Section
\ref{sec:method} (raw census of 18{,}198{,}978 reports / 3{,}740 outlets / newspapers
92.93\% retained; production target 18{,}198{,}978 articles, a pre-co-occurrence-filter
basis, 2007--2025, 110 banks).

\paragraph{Patent data.}
The innovation-output side is planned to link IncoPat and China National Intellectual
Property Administration patent data (enrichment M0.5.11): the primary outcome is
citation-weighted patents (\code{citation\_weighted\_patent}); patent-quality weighting
is reported on two parallel bases — \textcite{kogan2017technological} market-value
weighting for listed banks (requires stock prices, feasible only for listed banks) and
citation-count weighting in the Hall--Jaffe--Trajtenberg lineage for unlisted banks;
the secondary outcome is the fintech patent count (\code{patent\_fintech\_count}).

\paragraph{Financial and mediator data.}
The scenario-track paired outcome is noninterest income share
(\code{noninterest\_income\_share}) from the financial panel; the mediators are IT
capital expenditure (\code{capex\_it}) from CSMAR and R\&D intensity
(\code{r\_and\_d\_intensity}); the listed-bank extension (5-day CAR) requires CSMAR/Wind
daily returns and is not part of the annual mainline. Enrichment of bank-level financial
controls (log assets, leverage, capital ratio, NPL ratio, deposit growth) and year-level
macro controls (GDP growth, M2 growth) is in progress; the replication package
self-checks with \code{capture confirm} at the data-preparation step and assembles only
the columns that actually exist, so the controls are incorporated automatically once the
financial panel arrives.

\paragraph{Descriptive statistics.}
\NoteThis{Table 1 descriptive statistics (means / SDs / quantiles / bank coverage) to be
backfilled once the production-corpus (T=19) panel is generated.}

\section{Results}\label{sec:results}

\subsection{Main regressions: layer $\times$ event study (Table 2)}

Table 2 reports estimates of equations \eqref{eq:baseline} and \eqref{eq:eventstudy}:
the left panel presents the two-way fixed-effects baseline (columns 1--2:
\code{citation\_weighted\_patent}; columns 3--4: \code{noninterest\_income\_share}),
and the right panel the four-event layer $\times$ event studies. Reading frame: H1
predicts $\beta_1>0$ in columns 1--2; H2 predicts $\beta_2>0$ in columns 3--4 with
$\beta_2^{(patent)}$ markedly weaker than $\beta_1^{(patent)}$ (the placebo side of the
structural pairing); the event studies predict differential dynamics for
technology-layer-exposed banks after E1 (and after E4/E5a/E5b depending on regulatory
direction). \NoteThis{Table 2 coefficients, standard errors, joint pre-trend F, and
rwolf-adjusted p-values await production regressions.}

\subsection{Mediation decomposition (Table 3)}

Table 3 reports the Baron--Kenny three-step association decomposition of
\code{locus\_layer\_full} $\to$ \code{capex\_it} (primary mediator) /
\code{r\_and\_d\_intensity} (secondary mediator) $\to$
\code{citation\_weighted\_patent} \parencite{baron1986moderator}, with the
\textcite{imai2010general} sensitivity parameter $\rho$ (the threshold at which the
indirect effect changes sign under sequential-ignorability violations). The table is
explicitly labeled an exploratory association decomposition, not causal mediation.
\NoteThis{Table 3 direct/indirect decomposition and $\rho$ thresholds to be backfilled.}

\subsection{Heterogeneity: three waves and ownership (Table 4)}

Table 4 reports H4 by wave splits (pre 2007--2012 / wave\_1 2013--2015 / wave\_2
2016--2018 / consolidation 2019--2021 / wave\_3 2022--2024) and ownership (state-owned
vs.\ non-state): we expect the patent effect of \code{locus\_layer\_full} to concentrate
in wave\_2 and wave\_3, the noninterest-income effect of \code{locus\_scenario\_full}
to concentrate in wave\_1, with the consolidation period as comparison. All splits are
labeled exploratory. \NoteThis{Table 4 subgroup coefficients and between-group tests to
be backfilled.}

\subsection{Alternative classifiers and alternative estimators (Table 5)}

Table 5 reports side by side: (i) ctx / ss track alternative classifiers
(\code{tech\_frac\_lc} / \code{tech\_frac\_ls} and their derived indicators), each track
just-identified on its own; (ii) cross-vendor per-model classifiers; (iii) the
\textcite{borusyak2024revisiting} imputation DiD and the
\textcite{callaway2021difference} group-time estimator; (iv) the sibling paper's
collab\_mode as an alternative-specification control (control only, not a parallel
treatment). \NoteThis{Table 5 arm-by-arm coefficients and the cross-arm consistency
narrative to be backfilled.}

\subsection{IV estimates (2SLS)}

Following equation \eqref{eq:iv}, the six just-identified
endogenous-variable--instrument pairs are reported column by column with first-stage F,
second-stage coefficients, and (when instruments are weak) Anderson--Rubin confidence
sets. \NoteThis{IV table coefficients and KP F statistics to be backfilled.}

\subsection{Cell decomposition (Table 6)}

Table 6 reports the $2^2$ (4-cell) locus $\times$ collab\_mode cross-construct
decomposition, testing that the locus effect is not a shadow of the collaboration
dimension. \NoteThis{Table 6 cell means and decomposition shares to be backfilled;
paper\_spec states both ``$2^2$ cross-tab cells'' and ``8-cell decomposition'' — this
draft standardizes on the $2^2$ basis, with the cell count to be verified against the
production implementation.}

\section{Conclusion}\label{sec:conclusion}

This paper operationalizes the technology $\times$ business-function distinction of the
digital-finance taxonomy \parencite{gomber2017cube} as a binary locus measure on an
19-year media-discourse panel of 110 Chinese banks, and tests a
binary-treatment-to-binary-outcome structural pairing: technology-layer discourse
$\to$ citation-weighted patents (H1, confirmatory) and scenario discourse $\to$
noninterest income share (H2, exploratory pairing). \NoteThis{Concluding summary of
results (support for H1--H4) awaits production regressions.}

Theoretically, the results will speak to whether the layered-modular architecture
\parencite{yoo2010layered} and the information-versus-communication dichotomy
\parencite{boot2021fintech} leave an observable fingerprint on the form of bank
innovation output; if the structural pairing holds, ``doing fintech'' is not a
homogeneous scalar, and ignoring a bank's technology-layer positioning systematically
mismeasures the innovation returns to digital transformation. For policy, the results
will offer a layered perspective on fintech regulation: shocks targeting application
scenarios (payments / lending / wealth management) and shocks targeting data or
underlying technology (such as the two data-export regulations) may act on entirely
different margins of bank innovation; for bank strategy, the results will map the
correspondence between technology-discourse positioning and the form of innovation
output. \NoteThis{Directional policy statements to be tightened after results are
confirmed.}

Three limitations must be acknowledged. First, annual frequency: the panel is at the
bank--year level, so within-year discourse--investment dynamics are invisible and the
temporal resolution of the event studies is limited. Second, endogeneity of media
coverage: media discourse is jointly generated by banks' disclosure strategies and
media selection; although the leave-one-out peer instruments and event studies mitigate
part of the concern, the discourse measure is not randomly assigned, and the mediation
analysis is only an association decomposition. Third, external validity: China is a
single and distinctive digital-finance market (concentrated regulation, platform
dominance, a high share of state-owned banks), so extrapolation to other institutional
environments requires caution. In addition, the locus machine label is a hybrid of LLM
voting and keyword adjustment; its validity evidence is delivered through the
pre-committed expert-annotation checks (Appendix \ref{app:validity}): the LLM vote label
agrees with the expert majority 94.0\% of the time on the full track but only about 85\%
on the ctx and ss tracks, so window tracks carry larger locus error.

\appendix

\section{Measurement-Validity Commitments}\label{app:validity}

This appendix pre-registers the validity-checking protocol for the locus measure
(protocol pre-registered on 2026-06-22; expert annotation completed on 2026-10-12).
All intervals below are bootstrap 95\% intervals from 2{,}000 resamples of articles.

\paragraph{Criterion validity.}
Three coders independently label locus for a stratified sample of 575 articles on blind
workbooks with no machine labels; a third party adjudicates the 4 items on which the two
coders who coded the item disagree; the expert majority is compared with the LLM
three-track vote labels (pipeline stage 63, Panel F). Inter-coder Krippendorff's $\alpha$
is 0.856 [0.810, 0.899] with 92.7\% raw agreement. Full-track accuracy is 94.0\%
[91.3\%, 96.4\%], Cohen's $\kappa = 0.879$, Gwet's AC1 $= 0.882$, and 93.7\% after
reweighting to the population; ctx and ss accuracy is 84.9\% and 84.1\%
($\kappa = 0.685$ and 0.670). Coders' backgrounds and compensation were not recorded. Disclosure: the locus machine label is a HYBRID construction (LLM majority
vote $+$ keyword confidence adjustment $+$ empty-vote fallback); its
criterion/discriminant $\kappa$ inherits a keyword circularity absent from the pure-LLM
dimensions (substance / collab), so locus validity numbers are not directly comparable
to the pure-LLM dimensions.

\paragraph{Inter-track reliability.}
Machine--machine agreement across the full / ctx / ss tracks is reported as Cohen's
$\kappa$ (stage 71, all retained articles): 0.703 for full--ctx, 0.682 for
full--ss, and 0.776 for ctx--ss.

\paragraph{Construct orthogonality and discriminant validity.}
Article-level construct orthogonality uses normalized mutual information (NMI) as the
main evidence (stage 72); auxiliary discriminant evidence is the non-redundancy check of
$|corr|$ between the focal share and adjacent measures such as language clarity and
numerical intensity (stage 73; within-single-LLM-method non-redundancy evidence, not
interpreted as directional $\kappa$). \NoteThis{NMI and $|corr|$ values to be
backfilled.}

\section{Network Roles: Exploratory Appendix}\label{app:network}

This appendix is exploratory, not part of the main regressions. First, a mediation
channel: we take betweenness centrality in the bank co-occurrence network
(\code{bank\_between\_cent}) as a candidate mediator and ask whether technology-layer
discourse is associated with patent output through a cross-group knowledge-brokerage
position (a structural-holes logic) — a channel (knowledge brokerage $\to$ innovation)
explicitly distinguished from the risk-contagion channel ($\to$ profitability) studied
in a sibling paper. Because the mediator is endogenous, we estimate within the
association-decomposition framework and report \textcite{imai2010general} $\rho$
sensitivity. Second, heterogeneity: locus effects are split by the within-year hub
decile indicator (\code{bank\_is\_hub}). \NoteThis{Network mediation and hub
heterogeneity estimates await the production panel.}

Two disclosures. First, criterion sensitivity: as of 2026-07, network metrics come in
three bank--bank co-occurrence criterion variants — same article (unsuffixed),
$\pm 1$-sentence window (\code{\_win1}), and strict same sentence (\code{\_sent}); the
three criteria are mutually robust, each just-identified on its own, forming a
specification-robustness family side by side with no family-wise error correction and
no pooled over-identification. Second, reproducibility: the current-flow betweenness
centrality (\code{bank\_cfbetween\_cent}), clustering coefficient
(\code{bank\_clustering}), and harmonic centrality (\code{bank\_harmonic\_cent}) column
families (including their criterion-suffix variants) are BLAS floating-point
nondeterministic; cross-machine reproduction verifies zero drift after excluding these
columns by exact name. The Freeman betweenness centrality used in this paper
(\code{bank\_between\_cent}) is deterministic and is not on the exclusion list. The
network missingness convention is zero-filling off-network banks while disclosing the
on-network count $N$.

\printbibliography

\end{document}
